\section{Noise-stable spacetime Recognition and experimental reconstruction}
\label{sec:noise-stable-spacetime}

The exact identification theorem in Section~\ref{sec:metric-identification} determines the metric and connection from exact propagation, scale, and transport data.  Real observations are finite and noisy.  Recognition Kernel theory supplies the quantitative extension: after diffeomorphism and gauge directions are removed, a uniformly valid decoder converts finite data error into a bound on the physical quotient geometry.

Let $u_0=(g_0,A_0)$ be a smooth background and let $\mathscr S$ be a local Hilbert slice transverse to the diffeomorphism--gauge orbit through $u_0$.  A tangent vector in $T_u\mathscr S$ therefore represents a physical metric--connection perturbation rather than a pure lawful symmetry.  Let $\mathscr D$ be a Hilbert data space containing finite-resolution encodings of characteristic propagation, pointwise scale, and parallel transport, and let
\[
\cO_{\mathrm{geom}}:\mathscr S\longrightarrow\mathscr D
\]
be the corresponding observation map.

\begin{theorem}[Fixed-decoder spacetime Recognition theorem]
\label{thm:fixed-decoder-spacetime}
Let $B\subset\mathscr S$ be convex and suppose that $\cO_{\mathrm{geom}}$ is continuously Fr\'{e}chet differentiable on $B$.  Assume that there is a bounded linear decoder
\[
L:\mathscr D\longrightarrow\mathscr S
\]
and a number $0\le\delta<1$ such that, after identifying tangent spaces with the model Hilbert space of the slice,
\begin{equation}
\sup_{u\in B}
\norm{\id_{\mathscr S}-L D\cO_{\mathrm{geom}}(u)}
\le\delta.
\label{eq:uniform-decoder-defect}
\end{equation}
Then $\cO_{\mathrm{geom}}$ is injective on $B$ and, for all $u_1,u_2\in B$,
\begin{equation}
\norm{u_1-u_2}_{\mathscr S}
\le
\frac{\norm L}{1-\delta}
\norm{\cO_{\mathrm{geom}}(u_1)-\cO_{\mathrm{geom}}(u_2)}_{\mathscr D}.
\label{eq:fixed-decoder-stability}
\end{equation}
Thus a single decoder that remains an approximate left inverse throughout the neighbourhood gives a quantitative target-faithful reconstruction theorem.
\end{theorem}

\begin{proof}
Set $v=u_1-u_2$ and $u_t=u_2+tv$.  Convexity gives $u_t\in B$.  The fundamental theorem of calculus in Banach spaces yields
\[
\cO_{\mathrm{geom}}(u_1)-\cO_{\mathrm{geom}}(u_2)
=
\int_0^1D\cO_{\mathrm{geom}}(u_t)v\,\dd t.
\]
Applying $L$ and subtracting $v$ gives
\[
L\bigl(\cO_{\mathrm{geom}}(u_1)-\cO_{\mathrm{geom}}(u_2)\bigr)-v
=
\int_0^1
\bigl(LD\cO_{\mathrm{geom}}(u_t)-\id\bigr)v\,\dd t.
\]
Therefore \eqref{eq:uniform-decoder-defect} implies
\[
\norm{
L\bigl(\cO_{\mathrm{geom}}(u_1)-\cO_{\mathrm{geom}}(u_2)\bigr)-v
}
\le\delta\norm v.
\]
The reverse triangle inequality gives
\[
(1-\delta)\norm v
\le
\norm L\,
\norm{\cO_{\mathrm{geom}}(u_1)-\cO_{\mathrm{geom}}(u_2)},
\]
which is \eqref{eq:fixed-decoder-stability}.  If the two observations agree, the right-hand side vanishes and $u_1=u_2$.
\end{proof}

\begin{corollary}[Background decoder and small-ball reconstruction]
\label{cor:background-decoder-spacetime}
Let
\[
T_0:=D\cO_{\mathrm{geom}}(u_0):\mathscr S\to\mathscr D
\]
be bounded below:
\begin{equation}
\norm{T_0h}_{\mathscr D}
\ge\alpha_0\norm h_{\mathscr S},
\qquad \alpha_0>0.
\label{eq:background-observation-lower-bound}
\end{equation}
Then $T_0^*T_0\ge\alpha_0^2\id$ and the minimum-norm left inverse
\begin{equation}
L_0:=(T_0^*T_0)^{-1}T_0^*
\label{eq:background-minimum-decoder}
\end{equation}
is bounded with
\begin{equation}
L_0T_0=\id,
\qquad
\norm{L_0}\le\frac1{\alpha_0}.
\label{eq:background-decoder-norm}
\end{equation}
Suppose that on a convex neighbourhood $B$ one has
\begin{equation}
\sup_{u\in B}
\norm{D\cO_{\mathrm{geom}}(u)-T_0}
\le\eta
<\alpha_0.
\label{eq:derivative-closeness}
\end{equation}
Then
\begin{equation}
\norm{u_1-u_2}_{\mathscr S}
\le
\frac1{\alpha_0-\eta}
\norm{\cO_{\mathrm{geom}}(u_1)-\cO_{\mathrm{geom}}(u_2)}_{\mathscr D}
\label{eq:small-ball-geometric-stability}
\end{equation}
for all $u_1,u_2\in B$.
\end{corollary}

\begin{proof}
Inequality \eqref{eq:background-observation-lower-bound} implies
\[
\langle T_0^*T_0h,h\rangle
=\norm{T_0h}^2
\ge\alpha_0^2\norm h^2,
\]
so $T_0^*T_0$ is boundedly invertible and \eqref{eq:background-minimum-decoder} is well defined.  It is the Moore--Penrose left inverse on $\Ran T_0$, extended by zero on $(\Ran T_0)^\perp$, and its norm is at most $1/\alpha_0$.

For $u\in B$,
\[
\norm{\id-L_0D\cO_{\mathrm{geom}}(u)}
=
\norm{L_0(T_0-D\cO_{\mathrm{geom}}(u))}
\le\frac{\eta}{\alpha_0}.
\]
Thus Theorem~\ref{thm:fixed-decoder-spacetime} applies with
$\delta=\eta/\alpha_0$ and $\norm{L_0}\le1/\alpha_0$, giving
\[
\frac{\norm{L_0}}{1-\delta}
\le
\frac{1/\alpha_0}{1-\eta/\alpha_0}
=
\frac1{\alpha_0-\eta}.
\]
\end{proof}

\begin{corollary}[Lipschitz derivative criterion]
\label{cor:lipschitz-derivative-spacetime}
If
\begin{equation}
\norm{D\cO_{\mathrm{geom}}(u)-D\cO_{\mathrm{geom}}(v)}
\le L_D\norm{u-v}_{\mathscr S}
\label{eq:geometric-derivative-lipschitz}
\end{equation}
near $u_0$, then on every convex ball $B_r(u_0)$ with
\[
L_Dr<\alpha_0
\]
one may take $\eta=L_Dr$ in Corollary~\ref{cor:background-decoder-spacetime}.  Hence
\begin{equation}
\norm{u_1-u_2}_{\mathscr S}
\le
\frac1{\alpha_0-L_Dr}
\norm{\cO_{\mathrm{geom}}(u_1)-\cO_{\mathrm{geom}}(u_2)}_{\mathscr D}.
\label{eq:lipschitz-small-ball-stability}
\end{equation}
\end{corollary}

\begin{corollary}[Finite-error geometric reconstruction]
\label{cor:finite-error-spacetime}
Under the hypotheses of Corollary~\ref{cor:background-decoder-spacetime}, suppose measured data $d_i\in\mathscr D$ satisfy
\[
\norm{d_i-\cO_{\mathrm{geom}}(u_i)}_{\mathscr D}
\le\varepsilon_i.
\]
Then
\begin{equation}
\norm{u_1-u_2}_{\mathscr S}
\le
\frac{
\norm{d_1-d_2}_{\mathscr D}
+\varepsilon_1+\varepsilon_2
}{\alpha_0-\eta}.
\label{eq:noisy-geometric-reconstruction}
\end{equation}
In particular, if two models fit the same measured data to accuracy $\varepsilon$, then
\begin{equation}
\norm{u_1-u_2}_{\mathscr S}
\le
\frac{2\varepsilon}{\alpha_0-\eta}.
\label{eq:same-data-geometric-bound}
\end{equation}
\end{corollary}

\begin{proof}
The data errors give
\[
\norm{\cO_{\mathrm{geom}}(u_1)-\cO_{\mathrm{geom}}(u_2)}
\le
\norm{d_1-d_2}+\varepsilon_1+\varepsilon_2.
\]
Insert this estimate into \eqref{eq:small-ball-geometric-stability}.  The final statement is the case $d_1=d_2$ and $\varepsilon_1=\varepsilon_2=\varepsilon$.
\end{proof}

\begin{corollary}[Experimental Recognition Kernel criterion]
At a background for which
\[
\Ker D\cO_{\mathrm{geom}}(u_0)
=
\text{the diffeomorphism--gauge tangent space}
\]
and the induced observation on a transverse slice satisfies the positive lower bound \eqref{eq:background-observation-lower-bound}, the geometric target is locally target-faithful and stably decodable.  A residual larger than the calibrated right-hand side of \eqref{eq:noisy-geometric-reconstruction} rejects at least one declared propagation, scale, transport, model-class, gauge-slice, or uncertainty assumption.
\end{corollary}

\begin{remark}[Experimental boundary]
The theorem supplies a finite-error reconstruction contract and identifies the stability constants $\alpha_0$ and $\eta$ that an experiment or numerical protocol must estimate.  It does not claim that a particular instrument already measures complete characteristic or holonomy data, nor that the transverse lower bound is automatically positive.  Those are empirical and protocol-design obligations, now expressed as testable quantities rather than perfect-information prose.
\end{remark}
